\documentclass[12pt,a4paper,UTF8]{ctexart}
\usepackage[a4paper,top=22mm,bottom=22mm,left=22mm,right=22mm,headheight=15pt]{geometry}
\newcommand{\Needspace}[1]{\par\begingroup\dimen0=#1\relax\dimen2=\pagegoal\advance\dimen2 by-\pagetotal\ifdim\dimen0>\dimen2\newpage\fi\endgroup}
\newcommand{\needspace}[1]{\par\begingroup\skip0=\lastskip\vskip-\skip0\dimen0=#1\relax\advance\dimen0 by\skip0\vskip\dimen0\penalty-100\vskip-\dimen0\vskip\skip0\penalty9999\endgroup}
\usepackage{amsmath,amssymb,mathtools,bm}
\usepackage{graphicx,xcolor,booktabs,longtable,array,enumitem,fancyhdr,caption,float,microtype,setspace}
\usepackage[hidelinks,unicode]{hyperref}
\usepackage[most]{tcolorbox}
\definecolor{ContestBlue}{HTML}{17365D}
\hypersetup{pdftitle={第061套},pdfauthor={}}
\setstretch{1.24}
\setlength{\parindent}{2em}
\setlength{\parskip}{0.35em}
\setlength{\emergencystretch}{2em}
\setlist{itemsep=0.35em,topsep=0.4em}
\captionsetup{font=small,labelfont=bf,justification=centering}
\setcounter{secnumdepth}{0}
\pagestyle{fancy}
\fancyhf{}
\fancyhead[L]{\small 第061套}
\fancyfoot[C]{\thepage}
\renewcommand{\headrulewidth}{0.4pt}
\ctexset{
  section={format=\Large\bfseries\color{ContestBlue},beforeskip=1.2em,afterskip=0.7em},
  subsection={format=\large\bfseries\color{ContestBlue},beforeskip=1.0em,afterskip=0.5em},
  subsubsection={format=\normalsize\bfseries,beforeskip=0.8em,afterskip=0.35em}
}
\usepackage{calc}
\newcounter{none}
\let\originalsubsection\subsection
\renewcommand{\subsection}{\needspace{6\baselineskip}\originalsubsection}
\makeatletter
\newcommand{\tightlist}{\setlength{\itemsep}{0pt}\setlength{\parskip}{0pt}}
\newcommand{\maxwidth}{\ifdim\Gin@nat@width>\linewidth\linewidth\else\Gin@nat@width\fi}
\newcommand{\maxheight}{\ifdim\Gin@nat@height>0.72\textheight 0.72\textheight\else\Gin@nat@height\fi}
\newcommand{\pandocbounded}[1]{\begingroup\setkeys{Gin}{width=\maxwidth,height=\maxheight,keepaspectratio}#1\endgroup}
\makeatother
\begin{document}
\begin{titlepage}
  \centering
  \vspace*{58mm}
  {\Huge\bfseries\color{ContestBlue} 第061套\par}
  \vfill
  \vspace*{28mm}
\end{titlepage}
\section{第 1 题}\label{ux7b2c-1-ux9898}

一根半径为 \(r\)、质量为 \(m_1\) 的均匀实心圆轴，两侧同轴固定两个半径为
\(R\)、质量均为 \(m_2\)
的均匀实心圆盘，组成悠悠球。细绳绕在圆轴上，绳端悬挂质量为 \(m\)
的配重。悠悠球置于倾角为 \(\theta\) 的粗糙斜面上，细绳从轴的侧面引出。

（1）当 \(\theta=\theta_0\) 时，系统能够静止，求 \(\theta_0\)。

（2）求悠悠球与斜面间静摩擦因数的最小值。

（3）将斜面调为水平，用手保持系统静止后释放，求悠悠球轴心的初加速度
\(a_0\)。

（4）取
\(0<\theta_1<\theta_0\)，且摩擦足以保证纯滚动。系统稳定运动后，求悠悠球轴心沿斜面向上的加速度
\(a_1\)。

（5）求稳定运动时细绳与竖直方向的夹角 \(\alpha\)。

（6）求配重的加速度 \(\bm a_2\)。

\begin{figure}[H]
\centering
\includegraphics[width=0.86\linewidth,height=\textheight,keepaspectratio,alt={图1-1　复合悠悠球、斜面与悬挂配重的几何关系}]{../figures/question-01.pdf}
\caption*{图1-1　复合悠悠球、斜面与悬挂配重的几何关系}
\end{figure}

\newpage

\section{第 2 题}\label{ux7b2c-2-ux9898}

（1）六个等质量弹性小球依次串在一根无限长的竖直光滑杆上。某时刻解除约束，并可分别赋予任意大小、任意方向的初速度。各次碰撞均为一维完全弹性碰撞，求小球之间可能发生的最多碰撞次数。

（2）\(N\) 个等质量质点沿竖直方向等间距 \(h\) 排列，最低质点距地面也为
\(h\)。最上方质点从静止开始下落；每次碰撞均完全非弹性，碰后的质点团继续下落。求最终质点团撞地前的速度。取
\(N=6\) 给出结果。

（3）取 \(N=100\)，求相同过程的撞地速度。

\begin{figure}[H]
\centering
\includegraphics[width=0.752\linewidth,height=\textheight,keepaspectratio,alt={图2-1　六球碰撞链与等间距质点的逐级黏连坠落}]{../figures/question-02.pdf}
\caption*{图2-1　六球碰撞链与等间距质点的逐级黏连坠落}
\end{figure}

\newpage

\section{第 3 题}\label{ux7b2c-3-ux9898}

一颗彗星绕太阳运动，周期为 \(1200\,\mathrm y\)，近日点距太阳
\(r_1=0.1\,\mathrm{AU}\)。忽略除太阳外其他天体对彗星轨道的作用，并把地球公转轨道视为半径
\(1\,\mathrm{AU}\) 的圆。

（1）求彗星远日点距离 \(r_2\)。

（2）当地球轨道与彗星轨道相交时，求两轨道切线的锐角 \(\theta\)。

（3）彗星接近太阳并在地球公转方向前方、距地球 \(r_3\ll1\,\mathrm{AU}\)
时，仍以日心轨道速度运动。求太阳参考系和地球参考系中看到的彗星速度，单位取
\(\mathrm{AU/y}\)。

（4）彗星位于地球公转速度正前方，距地球 \(r_3=1.0\times10^8\,\mathrm m\)
时爆炸为质量比 \(9:1\)
的两块。爆炸使大小两块沿彗星日心轨道的瞬时切线分别向前、向后分离。忽略太阳对随后地心运动的引力，求使小块撞击地球所需释放的最小能量。取
\(g=10\,\mathrm{m/s^2}\)，\(R_E=6.4\times10^6\,\mathrm m\)，彗星总质量
\(m=1.0\times10^6\,\mathrm{kg}\)，\(1\,\mathrm{AU}=1.5\times10^{11}\,\mathrm m\)，一年
\(365\) 天。

\begin{figure}[H]
\centering
\includegraphics[width=0.867\linewidth,height=\textheight,keepaspectratio,alt={图3-1　高偏心率彗星轨道与地球圆轨道的位置、速度关系}]{../figures/question-03.pdf}
\caption*{图3-1　高偏心率彗星轨道与地球圆轨道的位置、速度关系}
\end{figure}

\newpage

\section{第 4 题}\label{ux7b2c-4-ux9898}

定义微元光束的集光度

\[
\mathrm dG=f(n)\,\mathrm dS_\perp\,\mathrm d\Omega,
\]

其中 \(n\) 为介质折射率，\(\mathrm dS_\perp\)
为面元在垂直光线方向上的投影，\(\mathrm d\Omega\) 为光线束张开的立体角。

（1）证明在折射率为 \(n\) 的均匀各向同性介质中，面元 \(\mathrm dS_1\)
向各方向均匀发出的光被面元 \(\mathrm dS_2\)
接收时，\(\mathrm dG_1=\mathrm dG_2\)。

（2）证明理想镜面反射过程中集光度守恒。

（3）证明从折射率 \(n_1\) 的介质折射到 \(n_2\)
的介质时，集光度守恒，并求 \(f(n)\)，无量纲比例系数取 \(1\)。

（4）由热力学第二定律说明集光度不能在无源理想光学系统中减小。为简化讨论可取真空
\(n=1\)。

\begin{figure}[H]
\centering
\includegraphics[width=0.9\linewidth,height=\textheight,keepaspectratio,alt={图4-1　面元间光束传递与介质界面的折射几何}]{../figures/question-04.pdf}
\caption*{图4-1　面元间光束传递与介质界面的折射几何}
\end{figure}

\newpage

\section{第 5 题}\label{ux7b2c-5-ux9898}

质量为 \(M\)、磁偶极矩大小为 \(m\) 的小磁铁，磁矩始终沿 \(x\)
轴方向。它从远处以速度 \(v_0\) 沿轴线射向位于 \(x=0\)
平面的圆形导线环。线环半径为 \(r\)、自感为
\(L\)、电阻忽略，初始电流为零。感应电流产生的磁场排斥小磁铁，使其减速。

（1）小磁铁位于轴线上、距线环中心为 \(x\) 时，求它在线环中产生的磁通量
\(\Phi(x)\)。

（2）求线环在小磁铁处产生的磁场 \(B(x)\)。

（3）求小磁铁所受轴向力 \(F(x)\)。

（4）求小磁铁恰好不能穿过线环的临界初速度 \(v_m\)。

\begin{figure}[H]
\centering
\includegraphics[width=0.84\linewidth,height=\textheight,keepaspectratio,alt={图5-1　轴向磁偶极接近无电阻自感线环}]{../figures/question-05.pdf}
\caption*{图5-1　轴向磁偶极接近无电阻自感线环}
\end{figure}

\newpage

\section{第 6 题}\label{ux7b2c-6-ux9898}

圆偏振光可用于分时式 3D
显示。约定迎着光传播方向观察时，电场矢量逆时针旋转为左旋圆偏振光。双折射波片的普通光、非普通光折射率分别为
\(n_o,n_e\)，且 \(n_o>n_e\)。

（1）一束左旋圆偏振光垂直入射理想平面镜，判断反射光的旋向并说明理由。

（2）一束经 \(y\)
方向偏振片后的线偏振光通过波片。求产生左旋圆偏振光所需的最小波片厚度
\(t\)，并说明光轴相对 \(y\) 轴的放置方向。

（3）实际系统用两个相互垂直放置的电控双折射晶体交替产生左右旋圆偏振。设设计波长处平均折射率为
\(n_e,n_o\)，不同颜色造成的偏差为
\(\Delta n_e,\Delta n_o\)，且偏差远小于
\(|n_o-n_e|\)。求最坏情形下错误旋向``鬼影''强度与正常信号强度之比。

\begin{figure}[H]
\centering
\includegraphics[width=0.9\linewidth,height=\textheight,keepaspectratio,alt={图6-1　偏振片、波片、电控双折射器件和检偏通道}]{../figures/question-06.pdf}
\caption*{图6-1　偏振片、波片、电控双折射器件和检偏通道}
\end{figure}

\newpage

\section{第 7 题}\label{ux7b2c-7-ux9898}

半径为 \(r\)
的无限长圆柱面上密绕导线，每一匝为一个斜平面与圆柱面的交线，构成闭合椭圆；椭圆长半轴为
\(r/\cos\theta\)，短半轴为 \(r\)。沿柱轴方向单位长度匝数为
\(n\)，导线电流为 \(I\)。

（1）求轴线上一点的轴向磁场 \(B_z\)。

（2）求柱内磁场 \((B_x,B_y,B_z)\) 中的 \(B_x\)。

（3）求柱内的 \(B_y\)。

（4）求柱外一点 \((\rho,\varphi)\)、\(\rho>r\) 处的横向磁场分量
\(B_\rho,B_\varphi\)，并说明 \(B_z\)。

\begin{figure}[H]
\centering
\includegraphics[width=0.9\linewidth,height=\textheight,keepaspectratio,alt={图7-1　圆柱面上的倾斜椭圆线匝及电流方向}]{../figures/question-07.pdf}
\caption*{图7-1　圆柱面上的倾斜椭圆线匝及电流方向}
\end{figure}

\newpage

\section{第 8 题}\label{ux7b2c-8-ux9898}

质量为 \(m_0\) 的质点在质量为 \(M\)
的史瓦西黑洞赤道面内运动。取黑洞中心为原点，设

\[
\Gamma=\left(1-\frac{2GM}{c^2r}\right)^{-1},
\qquad
E=\frac\gamma\Gamma m_0c^2,
\]

其中
\(\gamma=\mathrm dt/\mathrm d\tau\)，是题目采用的坐标时间率，区别于局域静止观察者测得的狭义相对论洛伦兹因子。已知能量---动量关系

\[
\Gamma E^2=m_0^2c^4+\Gamma P_r^2c^2+P_\varphi^2c^2,
\]

\(P_r,P_\varphi\) 采用狭义相对论形式。

（1）求黑洞事件视界半径 \(R_s\)。

（2）导出径向运动方程并写出有效势能 \(V_{\rm eff}(r)\)。

（3）求稳定圆轨道允许的最小半径 \(R_{\min}\)。

（4）质点在最小稳定圆轨道上沿切向向前发射频率为 \(\nu\)
的光子。求光子到达无穷远处的冲量参数 \(b\) 与静止接收者测得的频率
\(\nu'\)。

（5）简述质点此后的可能命运。

\begin{figure}[H]
\centering
\includegraphics[width=0.86\linewidth,height=\textheight,keepaspectratio,alt={图8-1　史瓦西黑洞赤道面内的圆轨道与切向发射}]{../figures/question-08.pdf}
\caption*{图8-1　史瓦西黑洞赤道面内的圆轨道与切向发射}
\end{figure}
\end{document}
